\documentclass[10pt,a4paper]{article}
\usepackage[margin=1.6cm,top=1.4cm,bottom=1.4cm]{geometry}
\usepackage{amsmath,amssymb}
\usepackage{graphicx}
\usepackage{hyperref}
\usepackage{microtype}
\usepackage{parskip}
\setlength{\parskip}{3pt}
\usepackage[numbers,sort&compress]{natbib}
\usepackage{titlesec}
\titlespacing{\section}{0pt}{4pt}{2pt}

\title{\vspace{-1.5em}\textbf{Q1: k-Means Clustering in High Dimensions}\\[2pt]
       \normalsize COL761 Assignment 2, 2026 \quad BITBLASTING}
\author{}\date{}

\begin{document}
\maketitle
\vspace{-2.5em}

% -------------------------------------------------------
\section*{(a) Selection of $k$}

For each $k \in \{1,\ldots,15\}$ we minimise the k-means objective (WCSS) via
Lloyd's algorithm with 10 random restarts and record the inertia
$I(k)=\sum_{j=1}^{k}\sum_{x_i\in C_j}\|x_i-\mu_j\|^2$.
Features are $z$-normalised before fitting to prevent scale dominance.

\textbf{Method.} The inertia curve is monotonically decreasing; the \emph{elbow}
marks where marginal gain drops sharply.  We formalise this via the
\emph{normalised acceleration ratio}:
\[
  r(k) = \frac{\Delta^2 \tilde{I}(k)}{\,|\tilde{I}(k)|\,+\varepsilon},
  \qquad
  \tilde{I}(k)=\frac{I(k)-I_{\min}}{I_{\max}-I_{\min}},
\]
where $\Delta^2$ denotes the second finite difference.  Dividing by
$|\tilde{I}(k)|$ normalises the curvature relative to the current inertia
magnitude, avoiding false positives in the steep, low-$k$ region.
The selected $k$ is $\operatorname{argmax}_k r(k)$.

The resulting elbow and the supplementary silhouette scores are shown in
Figure~\ref{fig:plot}.  Both criteria agree on the selected value.

\begin{figure}[h]
  \centering
  \includegraphics[width=0.90\linewidth,height=4.5cm,keepaspectratio]{plot.png}
  \caption{Left: WCSS vs.\ $k$ with the selected $k$ (red dashed line).
           Right: Silhouette score vs.\ $k$ (supplementary).}
  \label{fig:plot}
\end{figure}

% -------------------------------------------------------
\section*{(b) Structural and Statistical Assumptions}

\begin{enumerate}\setlength\itemsep{1pt}\setlength\topsep{1pt}
  \item \textbf{Spherical, isotropic clusters.} k-means minimises squared
        Euclidean distances, implicitly assuming clusters are convex and
        roughly equal-radius hyperballs.
  \item \textbf{Equal cluster sizes.} The Voronoi partition favours balanced
        clusters; small, dense clusters are absorbed by large ones.
  \item \textbf{Gaussian generative model.} WCSS is the MLE objective for a
        mixture of isotropic Gaussians $\mathcal{N}(\mu_j, \sigma^2 I)$ with
        equal mixing weights and shared variance.
  \item \textbf{Euclidean metric appropriateness.} In high dimensions,
        $\|x-\mu\|$ concentrates sharply around $\sqrt{d}\,\sigma$ for any
        $\mu$ (Lévy's lemma / measure concentration~\cite{blum2020foundations}),
        so inter- vs.\ intra-cluster distance ratios shrink, degrading
        separation.
  \item \textbf{$k$ known a priori.} The model has no mechanism to infer $k$
        from data; it must be supplied externally.
\end{enumerate}

% -------------------------------------------------------
\section*{(c) Limitations of the Model Selection Procedure}

\begin{itemize}\setlength\itemsep{1pt}\setlength\topsep{1pt}
  \item \textbf{Curse of dimensionality.} In high $d$, the volume of a ball of
        radius $r$ satisfies $V_d(r)\propto r^d/\Gamma(d/2+1)\to 0$ for fixed
        $r$ as $d\to\infty$~\cite{blum2020foundations}.  Distances concentrate,
        making the elbow extremely flat and hard to detect reliably.
  \item \textbf{Sensitivity to normalisation.} The elbow location can shift
        with feature scaling or normalisation choices.
  \item \textbf{Local optima.} Lloyd's algorithm is not guaranteed to find the
        global WCSS minimum; different restarts may yield different inertias,
        adding noise to the elbow curve.
  \item \textbf{No probabilistic model.} WCSS provides no likelihood, so formal
        model comparison (AIC/BIC) is unavailable without additional assumptions.
\end{itemize}

% -------------------------------------------------------
\section*{(d) Methodological Improvements and Alternatives}

\begin{itemize}\setlength\itemsep{1pt}\setlength\topsep{1pt}
  \item \textbf{BIC/AIC on GMM.} Replace k-means with a full Gaussian Mixture
        Model and select $k$ by BIC, which penalises model complexity
        explicitly~\cite{mclachlan2000mixture}.
  \item \textbf{Gap statistic.} Compare observed $\log I(k)$ against its
        expectation under a uniform null reference distribution; the optimal $k$
        is the smallest $k$ where the gap $G(k)=E[\log I_{\text{ref}}(k)]-\log
        I(k)$ exceeds $G(k+1)-s_{k+1}$~\cite{tibshirani2001estimating}.
  \item \textbf{Stability-based selection.} Cluster on random subsamples and
        measure clustering stability (e.g.\ adjusted Rand index); the most
        stable $k$ is preferred~\cite{ben2002stability}.
  \item \textbf{Dimensionality reduction pre-processing.} Applying PCA or
        UMAP before clustering mitigates distance concentration and often
        sharpens the elbow.
\end{itemize}

% -------------------------------------------------------
\section*{(e) Failure Modes of k-Means}

k-Means is inappropriate in the following well-known settings.

\textbf{Non-convex / manifold-shaped clusters} (e.g.\ two interlocking rings):
the Voronoi partition cannot recover non-convex cluster boundaries.
\emph{Detection:} fit a $k$-nearest-neighbour graph and check whether the
resulting connectivity components match the k-means labels; persistent
disagreement signals non-convexity.  \emph{Alternative:} spectral clustering
or DBSCAN.

\textbf{Clusters of very different sizes or densities}: large clusters
``steal'' points from small dense ones to lower WCSS.
\emph{Detection:} compare cluster cardinalities and intra-cluster variances; high
imbalance suggests misfit.  \emph{Alternative:} GMM with full covariance or
DBSCAN.

\textbf{High-dimensional data with distance concentration}: as $d$ grows, the
ratio $\max_i\|x-x_i\| / \min_i\|x-x_i\|\to 1$~\cite{beyer1999nearest}, so
nearest-centroid assignment becomes nearly random.
\emph{Detection:} compute the coefficient of variation of pairwise distances;
values near~0 signal concentration.  \emph{Mitigation:} use projected features
($d'\ll d$) or subspace clustering.

\textbf{Categorical or mixed-type features}: squared Euclidean distance is
semantically meaningless for discrete or ordinal variables.
\emph{Alternative:} k-modes or k-prototypes.

% -------------------------------------------------------
\bibliographystyle{plainnat}
\begin{thebibliography}{9}
\bibitem{blum2020foundations}
A.~Blum, J.~Hopcroft, R.~Kannan.
\textit{Foundations of Data Science}.
Cambridge University Press, 2020.

\bibitem{tibshirani2001estimating}
R.~Tibshirani, G.~Walther, T.~Hastie.
Estimating the number of clusters via the gap statistic.
\textit{JRSS-B}, 63(2):411--423, 2001.

\bibitem{mclachlan2000mixture}
G.~McLachlan, D.~Peel.
\textit{Finite Mixture Models}.
Wiley, 2000.

\bibitem{ben2002stability}
S.~Ben-David, U.~von Luxburg, D.~P{\'a}l.
A sober look at clustering stability.
\textit{COLT}, 2006.

\bibitem{beyer1999nearest}
K.~Beyer, J.~Goldstein, R.~Ramakrishnan, U.~Shaft.
When is ``nearest neighbor'' meaningful?
\textit{ICDT}, 1999.
\end{thebibliography}

\end{document}
